\documentclass[14pt]{extarticle}
\usepackage[T1]{fontenc}
\usepackage[margin=1in]{geometry}
\usepackage{amsthm,amsmath}
\usepackage{hyperref}

\newtheorem{theorem}{Theorem}[section]
\newtheorem{lemma}{Lemma}[section]
\newtheorem{definition}{Definition}[section]
\newtheorem*{example}{Example}
\newtheorem{exercise}{Exercise}[section]
\setcounter{section}{11}

\newcommand{\inv}[1]{#1^{-1}}
\newcommand{\join}[3][,]{#2_0 #1 #2_1 #1 \cdots #1 #2_{#3}}
\newcommand{\Times}[2]{\join[\times]{#1}{#2}}
\newcommand{\Oplus}[2]{\join[\oplus]{#1}{#2}}
\newcommand{\normalin}{\triangleleft}
\newcommand{\1}{\{e\}}
\newcommand{\set}[2]{\{ \ #1 \ | \ #2 \ \}}
\newcommand{\cyc}[1]{\langle #1 \rangle}

\DeclareMathOperator{\Abelian}{Abelian}
\DeclareMathOperator{\Inn}{Inn}
\DeclareMathOperator{\Aut}{Aut}
\DeclareMathOperator{\Ker}{Ker}
\DeclareMathOperator{\modu}{mod}
\DeclareMathOperator{\id}{id}
\DeclareMathOperator{\lcm}{lcm}

\begin{document}

\setcounter{exercise}{10}
\begin{exercise}
  Prove that any finite Abelian group $G$ can be expressed as the external direct product 
  of cyclic group of order $\join{n}{t - 1}$,
  where $n_{i - 1}$ divides $n_i$ for all $i \in [ 1 , t - 1 ]$.
\end{exercise}
\begin{proof}
  induction on $t$.
  \begin{itemize}
    \item Base: $t = 0$, trivial.
    \item Induction: Since $G$ can be (uniquely, up to isomorphism) expressed as the external direct product of
      cyclic groups of prime powered. Let $S$ be the set of those cyclic groups, $P$ and $Q$ are empty sets.
      First, take $H \in S$ which is isomorphic to $Z_{p^n}$ for some prime $p$ and postive $n$.
      % 将最大的，order 为 p 的幂的群加入 R.
      Then put the group which is isomorphic to $Z_{p^m}$ where $m$ is maximum to $P$,
      and move other groups which has form $Z_{p^i}$ from $S$ to $Q$.
      Repeat this procedure untile $S$ is empty.
      
      For any distinct $H \ K \in P$, $|H|$ is relative prime to $|K|$, since they 
      have form $Z_{p^i}$ with different $p$.
      Therefore, the product of elements in $P$ form a large cyclic group $R$.
      And the product of $Q$ form a finite Abelian group, by induction hypothesis,
      it can be expressed as $\join[\oplus]{K}{m}$, and they satisfy the property we want to prove.

      The last thing is proving $|K_0|$ divides $|R|$.
      We now consider the worse situation, $K_0$ is the product of $Q$,
      then elements of $Q$ are relative prime to each other, thus, for any prime $p$,
      $Q$ has at most one element of order power of $p$.
      For any $H \in Q$,
      $H$ has form $p^i$ and $H \in S$, then $|H|$ divides $|R|$ due to the
      way we construct $R$. So $|K_0|$ divides $|R|$.
      If $K_0$ is not the product of $Q$, then $|K_0|$ divides the order of the product of $Q$,
      therefore $|K_0|$ divides $|R|$.
  \end{itemize}
\end{proof}

\end{document}